\section{Задание 1. Разбиение IP-сети на подсети}

Требуется разбить исходную IP-сеть на 5 подсетей с заданным числом адресов в подсетях.

\vspace{1em}

\noindentИсходные данные для варианта 18:

\begin{itemize}
	\item исходная сеть: \mintinline{text}{194.3.192.0};
	\item маска сети: \mintinline{text}{255.255.255.0};
	\item по крайней мере 3 подсети должны включать по 32 адреса.
\end{itemize}

Запишем исходную маску сети в двоичном виде:

\begin{minted}{text}
11111111.11111111.11111111.00000000
\end{minted}

Первые 24 бита маски определяют адрес сети, а оставшиеся 8 бит отводятся
под адресацию внутри неё. Таким образом, исходная сеть содержит

\[
	2^8 = 256
\]

адресов. Разобьём это адресное пространство на одну подсеть, включающую
128 адресов, и четыре подсети, каждая из которых включает 32 адреса:

\[
	256 = 128 + 32 + 32 + 32 + 32.
\]

Такое разбиение образует 5 подсетей, полностью использует адресное
пространство исходной сети и удовлетворяет условию: не менее трёх подсетей
включают по 32 адреса.

Полученные подсети имеют следующий вид:

\begin{enumerate}
	\item Подсеть \mintinline{text}{194.3.192.0/25}:
	\begin{itemize}
		\item адреса узлов: \mintinline{text}{194.3.192.1--194.3.192.126};
		\item широковещательный адрес: \mintinline{text}{194.3.192.127};
		\item число доступных узлам адресов: $128 - 2 = 126$.
	\end{itemize}

	\item Подсеть \mintinline{text}{194.3.192.128/27}:
	\begin{itemize}
		\item адреса узлов: \mintinline{text}{194.3.192.129--194.3.192.158};
		\item широковещательный адрес: \mintinline{text}{194.3.192.159};
		\item число доступных узлам адресов: $32 - 2 = 30$.
	\end{itemize}

	\item Подсеть \mintinline{text}{194.3.192.160/27}:
	\begin{itemize}
		\item адреса узлов: \mintinline{text}{194.3.192.161--194.3.192.190};
		\item широковещательный адрес: \mintinline{text}{194.3.192.191};
		\item число доступных узлам адресов: $32 - 2 = 30$.
	\end{itemize}

	\item Подсеть \mintinline{text}{194.3.192.192/27}:
	\begin{itemize}
		\item адреса узлов: \mintinline{text}{194.3.192.193--194.3.192.222};
		\item широковещательный адрес: \mintinline{text}{194.3.192.223};
		\item число доступных узлам адресов: $32 - 2 = 30$.
	\end{itemize}

	\item Подсеть \mintinline{text}{194.3.192.224/27}:
	\begin{itemize}
		\item адреса узлов: \mintinline{text}{194.3.192.225--194.3.192.254};
		\item широковещательный адрес: \mintinline{text}{194.3.192.255};
		\item число доступных узлам адресов: $32 - 2 = 30$.
	\end{itemize}
\end{enumerate}

Результат разбиения представлен в табл.~\ref{tab:subnets}.

\begin{table}[H]
	\centering
	\caption{Результат разбиения исходной сети на подсети}
	\label{tab:subnets}
	\resizebox{\textwidth}{!}{%
	\begin{tblr}{
		colspec={Q[c,m] Q[l,m] Q[l,m] Q[c,m]},
		row{1}={font=\bfseries},
		hlines,
		vlines
	}
		Подсеть & Адрес & Маска & Число адресов \\
		1 & {\ttfamily 194.3.192.0}   & {\ttfamily 255.255.255.128} & 128 \\
		2 & {\ttfamily 194.3.192.128} & {\ttfamily 255.255.255.224} & 32 \\
		3 & {\ttfamily 194.3.192.160} & {\ttfamily 255.255.255.224} & 32 \\
		4 & {\ttfamily 194.3.192.192} & {\ttfamily 255.255.255.224} & 32 \\
		5 & {\ttfamily 194.3.192.224} & {\ttfamily 255.255.255.224} & 32
	\end{tblr}%
	}
\end{table}

Схема распределения адресного пространства приведена на
рис.~\ref{fig:subnet-allocation}. Высота каждого прямоугольника
пропорциональна полному числу адресов в соответствующей подсети.

\begin{figure}[H]
	\centering
	\resizebox{\textwidth}{!}{\begin{tikzpicture}[
	every node/.style={font=\small},
	subnet/.style={draw=black, line width=0.6pt, minimum width=11cm, inner sep=0pt},
	boundary/.style={font=\small\ttfamily, anchor=west}
]
	\node[subnet, fill=blue!10, minimum height=6.4cm, anchor=north west,
		align=center] (n1) at (0,0) {%
		\textbf{Подсеть 1}\\[0.3em]
		Адрес: {\ttfamily 194.3.192.0}\\
		Маска: {\ttfamily 255.255.255.128}
	};

	\node[subnet, fill=green!10, minimum height=1.6cm, anchor=north west,
		align=center, font=\scriptsize] (n2) at (0,-6.4) {%
		\textbf{Подсеть 2}\\
		Адрес: {\ttfamily 194.3.192.128}\\
		Маска: {\ttfamily 255.255.255.224}
	};
	\node[subnet, fill=blue!10, minimum height=1.6cm, anchor=north west,
		align=center, font=\scriptsize] (n3) at (0,-8.0) {%
		\textbf{Подсеть 3}\\
		Адрес: {\ttfamily 194.3.192.160}\\
		Маска: {\ttfamily 255.255.255.224}
	};
	\node[subnet, fill=green!10, minimum height=1.6cm, anchor=north west,
		align=center, font=\scriptsize] (n4) at (0,-9.6) {%
		\textbf{Подсеть 4}\\
		Адрес: {\ttfamily 194.3.192.192}\\
		Маска: {\ttfamily 255.255.255.224}
	};
	\node[subnet, fill=blue!10, minimum height=1.6cm, anchor=north west,
		align=center, font=\scriptsize] (n5) at (0,-11.2) {%
		\textbf{Подсеть 5}\\
		Адрес: {\ttfamily 194.3.192.224}\\
		Маска: {\ttfamily 255.255.255.224}
	};

	\node[boundary, anchor=north west] at (11.25,0) {194.3.192.0};
	\node[boundary, anchor=south west] at (11.25,-6.4) {194.3.192.127};
	\node[boundary, anchor=north west] at (11.25,-6.4) {194.3.192.128};
	\node[boundary, anchor=south west] at (11.25,-8.0) {194.3.192.159};
	\node[boundary, anchor=north west] at (11.25,-8.0) {194.3.192.160};
	\node[boundary, anchor=south west] at (11.25,-9.6) {194.3.192.191};
	\node[boundary, anchor=north west] at (11.25,-9.6) {194.3.192.192};
	\node[boundary, anchor=south west] at (11.25,-11.2) {194.3.192.223};
	\node[boundary, anchor=north west] at (11.25,-11.2) {194.3.192.224};
	\node[boundary, anchor=south west] at (11.25,-12.8) {194.3.192.255};
\end{tikzpicture}
}
	\caption{Распределение адресов исходной сети между подсетями}
	\label{fig:subnet-allocation}
\end{figure}
